% Build: pdflatex -interaction=nonstopmode -halt-on-error directed-theory.tex
% Run twice to resolve references. No external bibliography or images are needed.
\documentclass[11pt,a4paper]{article}
\usepackage[T1]{fontenc}
\usepackage{lmodern}
\usepackage[margin=25mm]{geometry}
\usepackage{amsmath,amssymb,mathtools}
\usepackage{microtype}
\usepackage{booktabs,tabularx,array}
\usepackage{enumitem}
\usepackage{listings}
\usepackage{tikz-cd}
\usepackage{needspace}
\usepackage{xurl}
\usepackage[colorlinks=true,linkcolor=blue!50!black,urlcolor=blue!50!black]{hyperref}
\usepackage{fancyhdr}
\pagestyle{fancy}
\fancyhf{}
\fancyhead[L]{\small Directed type theory in Arend}
\fancyhead[R]{\small Implementation rules}
\fancyfoot[C]{\thepage}
\setlength{\headheight}{14pt}
\setlength{\parindent}{0pt}
\setlength{\parskip}{5pt plus 1pt}
\setlength{\emergencystretch}{2em}
\setlist{nosep,leftmargin=*}
\lstset{basicstyle=\small\ttfamily,columns=fullflexible,keepspaces=true,
  showstringspaces=false,breaklines=true,frame=single,framerule=.25pt,
  xleftmargin=4pt,xrightmargin=4pt,aboveskip=8pt,belowskip=8pt}
\newcommand{\Cat}{\mathsf{Cat}}
\newcommand{\Type}{\mathsf{Type}}
\newcommand{\Set}{\mathsf{Set}}
\newcommand{\Prop}{\mathsf{Prop}}
\newcommand{\I}{\mathbb I}
\newcommand{\DI}{\mathbb D}
\newcommand{\dl}{0_{\!d}}
\newcommand{\dr}{1_{\!d}}
\newcommand{\il}{0_{\!I}}
\newcommand{\ir}{1_{\!I}}
\newcommand{\Path}{\mathsf{Path}}
\newcommand{\DPath}{\mathsf{DPath}}
\newcommand{\hpath}{\mathsf{path}}
\newcommand{\dpath}{\mathsf{dpath}}
\newcommand{\idd}{\mathsf{idd}}
\newcommand{\idp}{\mathsf{idp}}
\newcommand{\coe}{\mathsf{coe}}
\newcommand{\coep}{\mathsf{coe}^{+}}
\newcommand{\isop}{\mathsf{iso}^{+}}
\newcommand{\isocoe}{\mathsf{iso}^{+}_{\mathsf{coe}^{+}}}
\newcommand{\filltwo}{\mathsf{fill}_{2}}
\newcommand{\fillthree}{\mathsf{fill}_{3}}
\newcommand{\rezk}{\mathsf{rezk}}
\newcommand{\Hom}[3]{\mathsf{Hom}_{#1}(#2,#3)}
\newcommand{\arr}{\mathrel{\rightsquigarrow}}
\newcommand{\apd}{\mathbin{@_{d}}}
\newcommand{\lamplus}[1]{\lambda^{+}#1.\,}
\newcommand{\ty}{\ \mathsf{type}}
\newcommand{\ctx}{\ \mathsf{ctx}}
\newcommand{\FV}{\mathsf{FV}}
\newcommand{\rulelabel}[1]{\quad\text{\scriptsize(#1)}}
\newcommand{\infer}[3]{\frac{\begin{gathered}#1\end{gathered}}{#2}\rulelabel{#3}}
\newcommand{\code}[1]{\texttt{\detokenize{#1}}}
\newcommand{\source}[1]{\par{\footnotesize\raggedright\textit{Implementation:} #1\par}\par}
\newcommand{\blocktitle}[1]{\Needspace{5\baselineskip}\subsection*{#1}}
\title{\vspace{-12mm}Directed type theory in Arend\\[3pt]
  \large Formation, introduction, elimination, and computation rules}
\author{}
\date{Repository snapshot \texttt{a7a01117506d}\quad\textbullet\quad 22 September 2026}

\begin{document}
\maketitle
\vspace{-8mm}
\thispagestyle{plain}

This note extracts the directed core from the checked-out implementation of
Arend. It follows the rule-and-equation presentation of \S2.1 of
\href{https://arend-lang.github.io/assets/lang-paper.pdf}{\emph{Theorem prover
Arend}} by Fedor Part, Valery Isaev, and Sergey Sinchuk. The source of the
directed rules is this repository, especially \code{lib/Prelude.ard} and the
typechecker; the linked paper describes the earlier HoTT-I core.

The extension keeps the ordinary interval and equality and adds a directed
interval, morphisms, composition fillers, transport for families of groupoids,
directed univalence, and a Rezk principle. The rules below describe the
implemented primitives. Statements explicitly called \emph{derived} are
library constructions, not additional judgmental equalities.

\section{Contexts, universes, and variance}\label{sec:contexts}

\textbf{Judgments.} Write $\Gamma\vdash A\ty$ for a well-formed type,
$\Gamma\vdash a:A$ for typing, and $a\equiv b$ for judgmental equality.
In contrast, $a=_A b$ is the ordinary path type. As in the paper, ambient
contexts are omitted from equations. Substitution $A[i:=r]$ is abbreviated
to $A(r)$ when the bound variable is clear.

There are two kinds of variable: invariant $x:A$ and covariant $x:^{+}A$.
The latter may occur only in positions checked covariantly. This is a
syntactic restriction, even when $A$ happens to be a groupoid. There is no
contravariant binding annotation in this implementation.
Write $\Gamma^\circ$ for the available invariant context: entering an
invariant position hides the previously available covariant bindings.
New covariant binders introduced inside that position are still permitted.

\[
\infer{\Gamma^\circ\vdash A\ty}{\Gamma,x:A\ctx}{invariant extension}
\qquad
\infer{\Gamma\vdash A\ty}{\Gamma,x:^{+}A\ctx}{covariant extension}
\]
All contexts in these schemata are assumed well formed. In particular,
declaring an invariant parameter does not allow its type to depend on an
earlier covariant parameter.

\textbf{Universes.} $\Type_u$ is the universe of groupoidal types of size
$u$, and $\Cat_u$ permits categorical types. The relevant cumulative rules are
\[
\infer{\Gamma\vdash A:\Type_u}{\Gamma\vdash A:\Cat_u}{inclusion}
\qquad
\Gamma\vdash\Type_u:\Cat_{u+1},\qquad
\Gamma\vdash\Cat_u:\Cat_{u+1}.
\]
Both hierarchies are cumulative in $u$. The converse inclusion is absent;
in particular $\Type_u:\Type_{u+1}$ is rejected. The finite homotopy levels
remain: $\Prop\subseteq\Set_u\subseteq\Type_u$ and
$\Set_u:(1\text{-}\Type)_{u+1}$.
Unsubscripted $\Type$ and $\Cat$ below follow the implementation's unindexed
universe notation; they do not assert a particular finite size level.

\source{\code{BindingVariance}; \code{CheckTypeVisitor.clearCategoricalContext}
and \code{visitParameter}; \code{ConstLevel.succ}; \code{Sort.compare}.
Tests: \code{CatUniverseTest}, \code{CatPiTest}. Full paths are in Appendix~\ref{app:sources}.}

\blocktitle{Dependent functions and covariant sums}
Write $\Pi^{\nu}_{x:A}B$ for a dependent function type whose binder has mode
$\nu\in\{0,+\}$; $0$ is the ordinary invariant binder. Its introduction and
application have the following shapes, assuming the displayed function type
is well formed:
\[
\infer{\Gamma,x:^{\nu}A\vdash b:B}
 {\Gamma\vdash\lambda^{\nu}x.b:\Pi^{\nu}_{x:A}B}{function introduction}
\]
\[
\infer{\Gamma\vdash f:\Pi^{\nu}_{x:A}B\qquad
       \Gamma^{\nu}\vdash a:A}
 {\Gamma\vdash f\,a:B[x:=a]}{application}
\qquad \Gamma^{+}=\Gamma,\quad\Gamma^{0}=\Gamma^\circ.
\]
The usual function $\beta$ and $\eta$ equations apply, with variance respected.
A covariant function $f:A\to^{+}B$ therefore acts on morphisms by
\[
 f_{*}(p)\;\coloneqq\;\dpath(\lamplus{i}f(p\apd i)).
\]
An arbitrary invariant function $A\to B$ cannot be substituted for $f$ here.

There is an additional formation restriction: when a $\Pi$-type is checked
as a type expression, its domains are checked invariantly. Thus a covariantly
varying $A$ cannot in general be used to form $A\to^{+}B$ as an expression.
Declaration telescopes have a special checking rule for function-typed
parameters: a covariant parameter $f:^{+}A\to^{+}B$ may have a covariantly
varying $A$, as in the primitive $\isop$ below. This permission propagates
through successive covariant function domains, not through an invariant or
nested negative domain. It does not supply a general covariance rule for
function-type formation.

The covariant dependent sum is a categorical type former:
\[
\infer{\Gamma\vdash A:\Cat\qquad\Gamma,x:^{+}A\vdash B:\Cat}
 {\Gamma\vdash\Sigma^{+}_{x:A}B:\Cat}{covariant sum}
\]
\[
\infer{\Gamma\vdash a:A\qquad\Gamma\vdash b:B[a/x]}
 {\Gamma\vdash(a,b):\Sigma^{+}_{x:A}B}{pair}
\qquad
\infer{\Gamma\vdash s:\Sigma^{+}_{x:A}B}
 {\Gamma\vdash s.1:A\qquad\Gamma\vdash s.2:B[s.1/x]}{projections}
\]
\[
 (a,b).1\equiv a,\qquad (a,b).2\equiv b,\qquad(s.1,s.2)\equiv s.
\]
Arend writes this as \code{\Sigma+ (x : A) (B x)}: the mode is on the
sum, not on individual fields. Ordinary \code{\Sigma} uses invariant
fields and removes the categorical part of their universe levels; for
example, for invariant $C:\Cat$, $\Sigma_{x:C}\mathsf{Nat}:\Type$ whereas
$\Sigma^{+}_{x:C}\mathsf{Nat}:\Cat$. Covariant tuple elimination is
supported for $\Sigma^{+}$, but not for ordinary $\Sigma$.

\source{\code{CheckTypeVisitor.checkPi}, \code{visitTypeParameter},
\code{visitSigmaParameter}, \code{visitTuple};
\code{DefinitionTypechecker.typecheckParameters};
\code{StdImplicitArgsInference}. Tests: \code{CatPiTest}, \code{CatSigmaTest}.}

\section{The two intervals and their paths}\label{sec:paths}

\textbf{Ordinary paths.} The HoTT-I interval $\I$ (\code{I}) with endpoints
$\il$ (\code{left}) and $\ir$ (\code{right}) is retained. For an invariant
interval variable, $\Path(\lambda i.C,a,b)$ has introduction $\hpath$,
evaluation $p@i$, and the usual $\beta$, $\eta$, and endpoint equations.
The implementation allows the family to take values in $\Cat$.
Ordinary transport is
\[
\infer{\Gamma,i:\I\vdash C:\Cat\qquad
       \Gamma\vdash a:C(\il)\qquad\Gamma^\circ\vdash r:\I}
 {\Gamma\vdash\coe(\lambda i.C,a,r):C(r)}{ordinary transport}
\]
with $\coe(\lambda i.C,a,\il)\equiv a$ and
$\coe(\lambda i.C,a,r)\equiv a$ when $i\notin\FV(C)$.
This ordinary equality, rather than the directed morphism type, is used
for the homotopies in Sections~\ref{sec:univalence} and~\ref{sec:rezk}.

\textbf{Directed interval.} The new primitive is $\DI$ (\code{DI}), with
endpoints $\dl$ (\code{dleft}) and $\dr$ (\code{dright}):
\[
\infer{\Gamma\ctx}{\Gamma\vdash\DI:\Set_0}{directed interval}
\qquad
\infer{\Gamma\ctx}{\Gamma\vdash\dl:\DI}{left endpoint}
\qquad
\infer{\Gamma\ctx}{\Gamma\vdash\dr:\DI}{right endpoint}.
\]
A directed section uses a \emph{covariant} interval variable.
Given a covariant family $C$ and endpoints $a,b$, its morphism type has rules
\[
\infer{\Gamma,i:^{+}\DI\vdash C:\Cat\\
       \Gamma\vdash a:C(\dl)\qquad\Gamma\vdash b:C(\dr)}
 {\Gamma\vdash\DPath(\lamplus{i}C,a,b):\Type}{directed path formation}
\]
\[
\infer{\Gamma,i:^{+}\DI\vdash t:C}
 {\Gamma\vdash\dpath(\lamplus{i}t):
  \DPath(\lamplus{i}C,t(\dl),t(\dr))}{directed path introduction}
\]
\[
\infer{\Gamma\vdash p:\DPath(\lamplus{i}C,a,b)\qquad\Gamma\vdash r:\DI}
 {\Gamma\vdash p\apd r:C(r)}{directed evaluation}.
\]
The formation rule presupposes the well-formed family from its first
premise; the introduction rule has that same presupposition.
The defining equations are
\begin{align*}
 \dpath(\lamplus{i}t)\apd r&\equiv t[i:=r],
 &\dpath(\lamplus{i}p\apd i)&\equiv p\quad(i\notin\FV(p)),\\
 p\apd\dl&\equiv a,
 &p\apd\dr&\equiv b.
\end{align*}
For a constant family, put
\[
 \Hom{C}{a}{b}\;\coloneqq\;\DPath(\lamplus{i}C,a,b),\qquad
 \idd_a\;\coloneqq\;\dpath(\lamplus{i}a):\Hom{C}{a}{a}.
\]
We also write $a\arr b$ when $C$ is clear. Arend uses \code{a ~> b},
\code{idd}, and \code{p d@ i}; application \code{p i} elaborates to the
same directed evaluation. Below, $p(i)$ abbreviates $p\apd i$.
Morphisms in a category form a groupoidal type: the conclusion is $\Type$,
not an arbitrary $\Cat$.

\textbf{Elimination of the interval.} Matching on \code{DI} endpoints is
allowed for an invariant parameter $i:\DI$. With a well-formed invariant
motive $B$, branches $b_0:B(\dl)$ and $b_1:B(\dr)$ give
\[
 \mathsf{case}_{B}(\dl;b_0,b_1)\equiv b_0,
 \qquad\mathsf{case}_{B}(\dr;b_0,b_1)\equiv b_1.
\]
Such an eliminator takes its interval argument invariantly. Endpoint
matching on $i:^{+}\DI$ is rejected, so this is not an eliminator for
arbitrary directed sections. The partial boundary specifications of the
primitive fillers and transport below are installed specially by the
prelude implementation.

\source{\code{lib/Prelude.ard}: \code{DI}, \code{DPath}, \code{dpath},
\code{~>}, \code{idd}, \code{d@}; \code{AtExpression.make};
\code{CompareVisitor.comparePathEta}; \code{PatternTypechecking};
\code{ElimTypechecking}. Tests: \code{CatPreludeTest}.}

\clearpage
\section{Composition fillers}\label{sec:fillers}

\blocktitle{A two-dimensional filler}
For composable morphisms, there is a chosen square with three prescribed
edges:
\[
\infer{\Gamma\vdash C:\Cat\qquad\Gamma\vdash a,b,c:C\\
       \Gamma\vdash f:a\arr b\qquad\Gamma\vdash g:b\arr c\\
       \Gamma\vdash r,s:\DI}
 {\Gamma\vdash\filltwo(f,g,r,s):C}{\code{fill2}}
\]
\begin{align*}
 \filltwo(f,g,\dl,j)&\equiv a,\\
 \filltwo(f,g,\dr,j)&\equiv g(j),\\
 \filltwo(f,g,i,\dl)&\equiv f(i).
\end{align*}
All parameters in this primitive, including the category, objects, arrows,
and the two interval coordinates, are covariant. The square is oriented as
follows; its fourth edge is the composite.
\[
\begin{tikzcd}[row sep=large,column sep=huge]
 a \arrow[r,dashed,"g\circ f"] & c\\
 a \arrow[u,"\idd_a"] \arrow[r,"f"'] & b \arrow[u,"g"']
\end{tikzcd}
\qquad
\begin{gathered}
 i:\ \text{left to right},\\
 j:\ \text{bottom to top}.
\end{gathered}
\]
\textbf{Derived composition:}
\[
 g\circ f\;\coloneqq\;
 \dpath(\lamplus{i}\filltwo(f,g,i,\dr)):a\arr c.
\]
There is no further primitive reduction for the free edge $j=\dr$.
In particular, the unit and associativity laws are ordinary paths proved
in the library, not extra judgmental equations for $\filltwo$.

\blocktitle{A three-dimensional filler}
The primitive \code{fill3} fills a particular five-face boundary, with
one collapsed edge. Its exact boundary is needed: it is not an unrestricted
cubical horn-filling rule.

Fix $C:\Cat$ and seven objects $a,b,d,a',b',c',d':C$ with arrows
\[
\begin{array}{lll}
 ab:a\arr b,& bd:b\arr d,& ad:a\arr d,\\
 a'b':a'\arr b',& b'd':b'\arr d',& a'c':a'\arr c',\\
 c'd':c'\arr d',& aa':a\arr a',& bb':b\arr b',\\
 ac':a\arr c',& dd':d\arr d'.&
\end{array}
\]
These are names of arrows, not compositions. Supply the following faces:
\begin{align*}
 L&:\DPath(\lamplus{k}(a\arr a'c'(k)),aa',ac'),\\
 R&:\DPath(\lamplus{k}(bd(k)\arr b'd'(k)),bb',dd'),\\
 B&:\DPath(\lamplus{i}(ab(i)\arr ad(i)),\idd_a,bd),\\
 F&:\DPath(\lamplus{i}(a'b'(i)\arr c'd'(i)),a'c',b'd'),\\
 T&:\DPath(\lamplus{i}(ab(i)\arr a'b'(i)),aa',bb').
\end{align*}
Let $\partial$ denote exactly these objects, arrows, and five faces. Then
\[
\infer{\Gamma\vdash\partial\ \text{as above}\qquad
       \Gamma\vdash r,s,t:\DI}
 {\Gamma\vdash\fillthree(\partial;r,s,t):C}{\code{fill3}}
\]
has the judgmental boundary equations
\begin{align*}
 \fillthree(\partial;\dl,j,k)&\equiv L(k)(j),&
 \fillthree(\partial;\dr,j,k)&\equiv R(k)(j),\\
 \fillthree(\partial;i,\dl,k)&\equiv B(i)(k),&
 \fillthree(\partial;i,\dr,k)&\equiv F(i)(k),\\
 \fillthree(\partial;i,j,\dl)&\equiv T(i)(j).&&
\end{align*}
The face types enforce agreement on all overlaps. In particular,
$\fillthree(\partial;\dl,\dl,k)\equiv a$: this is the collapsed edge.
The missing $k=\dr$ face yields the derived term
\[
 \dpath\bigl(\lamplus{i}\dpath(\lamplus{j}
              \fillthree(\partial;i,j,\dr))\bigr)
 :\DPath(\lamplus{i}(ad(i)\arr c'd'(i)),ac',dd').
\]
All the mathematical parameters above are covariant. The concrete prelude
also has one invariant dummy argument of unit type, supplied as
\code{fill3 ()}; it is omitted from the mathematical notation.

\textbf{Derived coherence.} The library uses $\filltwo$ and $\fillthree$
to prove contractibility of the type of a composite together with a
composition triangle:
\[
 \Sigma^{+}_{h:a\arr c}
   \DPath(\lamplus{i}(f(i)\arr h(i)),\idd_a,g).
\]
Its chosen point is
$(g\circ f,\dpath(\lamplus{i}\dpath(\lamplus{j}\filltwo(f,g,i,j))))$.
Consequently the library proves
$f\circ\idd_a=f$, $\idd_b\circ f=f$, and
$h\circ(g\circ f)=(h\circ g)\circ f$ as ordinary equalities.

\source{\code{lib/Prelude.ard}: \code{fill2}, \code{fill3}; corresponding
cases of \code{Prelude.update}. Derived constructions:
\code{arend-lib/src/InfCat/Basics.ard}, especially \code{comp},
\code{segal2Prop}, \code{comp-char}, \code{comp_idd-left},
\code{comp_idd-right}, and \code{comp-assoc}.}

\section{Directed transport and directed univalence}\label{sec:univalence}

\blocktitle{Forward transport in a family of groupoids}
\[
\infer{\Gamma,i:^{+}\DI\vdash A:\Type_u\qquad
       \Gamma\vdash a:A(\dl)\qquad\Gamma\vdash r:\DI}
 {\Gamma\vdash\coep(\lamplus{i}A,a,r):A(r)}{\code{coe+}}
\]
\[
 \coep(\lamplus{i}A,a,\dl)\equiv a.
\]
The family, starting element, and interval argument are covariant.
The family must take values in $\Type_u$; there is no primitive of this
signature for arbitrary $A:\DI\to^{+}\Cat_u$, nor a primitive backward
transport along a general directed arrow.

\textbf{No general regularity equation.} Unlike ordinary $\coe$, the
normalizer does \emph{not} impose
$\coep(\lamplus{i}A,a,r)\equiv a$ for a constant family. Only the endpoint
rule above and the $\isop$ rule below are specified. This is an intentional
distinction in this description of the implementation, irrespective of
which propositional equalities can be derived.

\Needspace{16\baselineskip}
\blocktitle{A function determines a directed path of types}
\[
\infer{\Gamma\vdash A,B:\Type_u\qquad
       \Gamma\vdash f:A\to^{+}B\qquad\Gamma\vdash r:\DI}
 {\Gamma\vdash\isop(A,B,f,r):\Type_u}{\code{iso+}}
\]
This uses the covariant declaration-parameter convention described in
Section~\ref{sec:contexts}: $A,B,f,r$ are all covariant parameters.
Implicit $A,B$ are normally suppressed. Its equations are
\begin{align*}
 \isop(A,B,f,\dl)&\equiv A,&
 \isop(A,B,f,\dr)&\equiv B,\\
 \coep(\lamplus{i}\isop(A,B,f,i),a,\dr)&\equiv f(a)
 &&\bigl(i\notin\FV(A,B,f)\bigr).
\end{align*}
Only a forward function is required. No inverse or inverse homotopies
occur in this directed principle. Define
\[
 \mathsf{ua}^{+}(f)\coloneqq\dpath(\lamplus{i}\isop(A,B,f,i)):A\arr B,
 \qquad
 \mathsf{tr}(p)\coloneqq\lamplus{a}
       \coep(\lamplus{i}p(i),a,\dr).
\]
Then $\mathsf{tr}(\mathsf{ua}^{+}(f))\equiv f$ by the computation rule
and function $\eta$.

\blocktitle{The other composite is homotopic to the identity}
For $A,B:\Type_u$ and $p:A\arr B$, the primitive
\code{iso+_coe+} supplies a mixed directed/ordinary square:
\[
\infer{\Gamma\vdash A,B:\Type_u\qquad\Gamma\vdash p:A\arr B\\
       \Gamma\vdash r:\DI\qquad\Gamma^\circ\vdash s:\I}
 {\Gamma\vdash\isocoe(p,r,s):\Type_u}{\code{iso+_coe+}}
\]
The parameters $A,B,p,r$ are covariant, while $s$ is invariant. Its four
boundary equations are
\begin{align*}
 \isocoe(p,\dl,j)&\equiv A,&
 \isocoe(p,\dr,j)&\equiv B,\\
 \isocoe(p,i,\il)&\equiv\isop(A,B,\mathsf{tr}(p),i),&
 \isocoe(p,i,\ir)&\equiv p(i).
\end{align*}
Thus the other round trip is witnessed by an \emph{ordinary path}:
\[
 \hpath\bigl(\lambda j.\dpath(\lamplus{i}\isocoe(p,i,j))\bigr)
 :\mathsf{ua}^{+}(\mathsf{tr}(p))=p.
\]
In an invariant context, these give a function/path correspondence
$(A\to^{+}B)\simeq\Hom{\Type_u}{A}{B}$. This correspondence is a consequence
of the three primitives, with one judgmental round trip and one supplied
homotopy; it is not a judgmental $\eta$ rule for $\isop$.

\source{\code{lib/Prelude.ard}: \code{coe+}, \code{iso+}, \code{iso+_coe+};
\code{Prelude.update}; \code{NormalizeVisitor.normalizeCoe}, including its
exclusion of \code{COERCE_PLUS} from constant-family reduction.
Tests: \code{CatPreludeTest.coePlusTest}, \code{isoPlusTest},
\code{isoPlusLevelError}, \code{isoCoePlusTest}.}

\clearpage
\section{Rezk completeness and the groupoidal fragment}\label{sec:rezk}

\blocktitle{The Rezk principle}
Let $C:\Cat$, $x,y:C$, $f:x\arr y$, and $g:y\arr x$. Instead of merely
asking for equations between composites, the primitive takes two
composition triangles:
\[
 P:\DPath(\lamplus{i}(f(i)\arr x),\idd_x,g),\qquad
 Q:\DPath(\lamplus{i}(g(i)\arr y),\idd_y,f).
\]
Writing $S_x\coloneqq\Sigma^{+}_{z:C}\Hom{C}{x}{z}$, the rule is
\[
\infer{\Gamma\vdash C:\Cat\qquad\Gamma\vdash x,y:C\\
       \Gamma\vdash f:x\arr y\qquad\Gamma\vdash g:y\arr x\\
       \Gamma\vdash P,Q\ \text{of the displayed triangle types}}
 {\Gamma\vdash\rezk(f,g,P,Q):(x,\idd_x)=_{S_x}(y,f)}{\code{rezk}}.
\]
All arguments are covariant. Projecting this ordinary path to its first
component yields $x=_C y$; the second component identifies the chosen
arrow with the arrow represented by that equality. Thus the primitive
records more than a bare identification of objects.

\textbf{Computation.} No defining body or special judgmental computation
rule is installed for $\rezk$. Its ordinary path evaluations still satisfy
the general endpoint equations:
\[
 \rezk(f,g,P,Q)@\il\equiv(x,\idd_x),\qquad
 \rezk(f,g,P,Q)@\ir\equiv(y,f).
\]
Using the composition-triangle characterization of Section~\ref{sec:fillers},
the library derives that an invertible morphism is represented by an
ordinary equality. In particular, it proves the equivalence between
$x=_C y$ and the type of isomorphisms from $x$ to $y$.

\blocktitle{Identity elimination for morphisms in a groupoid}
When $A:\Type$, a morphism in a \emph{constant} family admits identity
elimination. A representative based form is
\[
\infer{\Gamma\vdash A:\Type\qquad\Gamma\vdash a:A\\
       \Gamma,y:A,p:a\arr y\vdash B(y,p)\ty\\
       \Gamma\vdash b:B(a,\idd_a)\qquad
       \Gamma\vdash y:A\qquad\Gamma\vdash p:a\arr y}
 {\Gamma\vdash J_d(B,b,y,p):B(y,p)}{groupoidal directed identity elimination}
\]
\[
 J_d(B,b,a,\idd_a)\equiv b.
\]
This schema is stated in the invariant fragment. In the syntax it is a
definition by the pattern \code{idd}. The checker requires a constant
ambient family whose type has a noncategorical sort and an endpoint
variable that can be eliminated, with the usual dependency/occurs checks.
It rejects the analogous general rule for $A:\Cat$, as well as an
unrestricted fixed-endpoint $K$ rule.

Ordinary identity elimination gives a map $(a=_C b)\to(a\arr b)$ for
any $C:\Cat$. For $A:\Type$, the rule above also gives the reverse map,
and the library proves
\[
 \Hom{A}{a}{b}\simeq(a=_A b).
\]
This explains the groupoidal fragment without identifying directed
morphisms with ordinary equality in a general category.

\source{\code{lib/Prelude.ard}: \code{rezk}; \code{Prelude.update};
\code{PatternTypechecking} and \code{CoreExpressionChecker} for
\code{idd} patterns. Tests: \code{IddTest}, especially \code{jTest},
\code{kTest}, \code{notInTypeError}, and \code{inTypeOk}.
Derived results: \code{InfCat/Basics.ard} and \code{InfCat/Rezk.ard}.}

\clearpage
\appendix
\section{Implementation qualifications and source map}\label{app:sources}

\textbf{Size of path types.} The formation rule in Section~\ref{sec:paths}
deliberately uses unindexed $\Type$. In an invariant context the size of
$\DPath(C,a,b)$ is computed from its family, with groupoidal homotopy
level. Unlike ordinary $\Path$, $\DPath$ does not lower a finite homotopy
level. If the path arguments depend on currently available covariant
variables, the elaborator forces both path forms to unindexed $\Type$
(infinite size and groupoidal homotopy level). It does not preserve the
family's finite size in that case. For example, invariant
$x:\mathsf{Nat}$ permits $x\arr x:\Type_0$, while covariant
$x:^{+}\mathsf{Nat}$ permits $x\arr x:\Type$ but not that finite-size
conclusion. This is implemented by \code{DefCallResult.isForcedInfinite}
and \code{GetTypeVisitor.visitPathType}, and tested by
\code{CatPreludeTest.dPathCovariantTest} and its variants.

\textbf{Scope.} This is a presentation of the directed core and the
structural rules it needs. Arend's general inductive definitions, records,
truncation, level proofs, termination, and elaboration have further
restrictions; their full rule systems are outside this note. For example,
covariant datatype parameters do not automatically permit constructors
with conditions or truncated declarations. The ordinary univalence
primitive \code{iso} also remains in the prelude, for groupoidal types.
No normalization, canonicity, or model-theoretic completeness theorem is
claimed here.

The following paths are relative to the repository root at commit
\code{a7a01117506d}. The full identifier is\par
{\small\code{a7a01117506d7c81c7d38b78feec5204af7ec949}.}\par
The working tree had no
tracked source modifications when inspected.

Java paths in the list below, unless prefixed by \code{api/}, are relative to
\path|base/src/main/java/org/arend/|.

\begingroup\footnotesize\raggedright
\begin{description}[style=nextline,font=\normalfont\bfseries,itemsep=4pt]
\item[Primitive signatures and boundaries]
\path|lib/Prelude.ard|;\newline
\path|prelude/Prelude.java|,
method \code{update} (the directed cases are adjacent).
\item[Variance and type formation]
\path|api/src/main/java/org/arend/ext/core/context/BindingVariance.java|;\newline
\path|typechecking/visitor/CheckTypeVisitor.java|;\newline
\path|typechecking/visitor/DefinitionTypechecker.java|;\newline
\path|typechecking/implicitargs/StdImplicitArgsInference.java|.
\item[Universe levels]
\path|api/src/main/java/org/arend/ext/core/level/ConstLevel.java|;\newline
\path|core/sort/Sort.java|;\newline
\path|core/sort/SortExpression.java|.
\item[Path typing, evaluation, and conversion]
\path|typechecking/result/DefCallResult.java|;\newline
\path|core/expr/AtExpression.java|;\newline
\path|core/expr/visitor/GetTypeVisitor.java|;\newline
\path|core/expr/visitor/CompareVisitor.java|;\newline
\path|core/expr/visitor/NormalizeVisitor.java|.
\item[Elimination restrictions]
\path|typechecking/patternmatching/PatternTypechecking.java|;\newline
\path|typechecking/patternmatching/ElimTypechecking.java|;\newline
\path|typechecking/doubleChecker/CoreExpressionChecker.java|.
\item[Regression tests]
\path|src/test/java/org/arend/cat/| (\code{CatPreludeTest}, \code{CatUniverseTest},
\code{CatPiTest}, \code{CatSigmaTest}, and related tests);\newline
\path|src/test/java/org/arend/typechecking/patternmatching/IddTest.java|.
\item[Derived categorical laws]
\path|arend-lib/src/InfCat/Basics.ard|;\newline
\path|arend-lib/src/InfCat/Rezk.ard|;\newline
\path|arend-lib/src/InfCat/Groupoid.ard|.
\end{description}
\endgroup

\clearpage
\section{Concrete syntax examples}\label{app:examples}

These examples use only the prelude. ASCII \code{:+}, \code{->+},
\code{=>+}, and \code{\Sigma+} are the counterparts of the superscript-plus
forms in \code{lib/Prelude.ard}. Expected types let the checker infer the
variance of the unannotated lambda binders in these examples.

\begin{lstlisting}
\func identity {C :+ \Cat} {a :+ C} : a ~> a
  => dpath (\lam _ => a)

\func map {C D :+ \Cat} (f :+ C ->+ D)
          {a b :+ C} (p :+ a ~> b) : f a ~> f b
  => dpath (\lam i => f (p i))

\func compose {C :+ \Cat} {a b c :+ C}
              (f :+ a ~> b) (g :+ b ~> c) : a ~> c
  => dpath (\lam i => fill2 f g i dright)

\func uaPlus.{u} {A B :+ \Type u} (f :+ A ->+ B) : A ~> B
  => dpath (\lam i => iso+ f i)

\func uaBeta.{u} {A B :+ \Type u}
                 (f :+ A ->+ B) (a :+ A)
  : coe+ (\lam i => iso+ f i) a dright = f a
  => idp

\func uaEta.{u} {A B :+ \Type u} (p :+ A ~> B)
  : dpath (\lam i =>
      iso+ (\lam a => coe+ (\lam k => p k) a dright) i) = p
  => path (\lam j => dpath (\lam i => iso+_coe+ p i j))

\func directedJ {A : \Type} {a : A}
  (B : \Pi {b : A} -> a ~> b -> \Type)
  (base : B idd) {b : A} (p : a ~> b) : B p \elim p
  | idd => base
\end{lstlisting}

The equations witnessed by \code{idp} in these examples are judgmental;
\code{uaEta} instead explicitly constructs an ordinary path.
The supplied PDF is compiled from this self-contained source. Rebuild
with two runs of \code{pdflatex directed-theory.tex}; see the accompanying
\code{README.md} for validation commands.

\clearpage
\section{A consistency counterexample}\label{app:inconsistency}

\textbf{The inspected snapshot is inconsistent.} The following prelude-only
module defines a closed term of a datatype with no constructors. It was
accepted by the repository's typechecker and its core double checker at
commit \code{a7a01117506d}. It is also saved as
\code{docs/directed/inconsistency.ard}.

\begin{lstlisting}
\data Empty

\func Family (i : DI) : \Type0 \elim i
  | dleft => \Sigma
  | dright => Empty

\func collapse {i : DI} (p : dleft ~> i) : Family i \elim p
  | idd => ()

\func contradiction : Empty
  => collapse (dpath (\lam i => i))
\end{lstlisting}

The derivation uses three features together:
\begin{enumerate}[itemsep=6pt]
\item The prelude declares \code{DI} as a datatype with two nullary
constructors. Its inferred sort is \code{\Set0}, and
\code{Prelude.update} does not replace that sort with a categorical one.
Thus $\DI:\Set_0\subseteq\Type$.
\item Directed path introduction gives the canonical segment
\[
 s\coloneqq\dpath(\lamplus{i}i):\dl\arr\dr.
\]
Since $\DI$ is classified as groupoidal, the \code{idd} pattern accepts
the based eliminator \code{collapse}.
\item Invariant endpoint matching gives
$\mathsf{Family}(\dl)\equiv\mathbf{1}$ and
$\mathsf{Family}(\dr)\equiv\mathbf{0}$. The identity case of
\code{collapse} is therefore the unit value $()$. Applying it to $s$
gives a term of $\mathsf{Family}(\dr)\equiv\mathbf{0}$.
\end{enumerate}

There are no user axioms, holes, recursive definitions, or additional
library imports in this module. Neither composition fillers nor directed
univalence nor the Rezk primitive is needed for the contradiction.

The decisive implementation locations are the \code{"DI"} case of
\code{Prelude.update} and the \code{typeSort.getHLevel().isCat()} check
for directed identity patterns in \code{PatternTypechecking}. A repair
must break this combination of rules; rejecting groupoidal elimination
on the directed interval would block this particular derivation.
That observation alone is not a consistency proof for a repaired system.

The 234 passing regression tests and the examples in Appendix~\ref{app:examples}
remain valid reports of checker behavior. They do not establish logical
consistency. This counterexample concerns the implemented snapshot; a
revised directed theory requires a separate soundness argument.

\end{document}
